\section{Verification}
\label{sec:verification}

\subsection{Unit Tests}
The arithmetic of the mission is kept in two modules that do not import ROS~2: the line
mathematics (offset, marker detection and steering) and the table of stations and colours.
They are covered by 13 unit tests, which run in under a second without the simulator. The
tests fix the behaviour that the rest of the mission relies on:
\begin{itemize}
  \item a floor without a line gives no offset, and a line between the middle channels gives
        an offset of zero;
  \item a line under the leftmost channel gives a positive offset and one under the
        rightmost a negative offset, which fixes the sign convention against the channel
        order of the robot description;
  \item readings below the floor threshold are ignored;
  \item a marker requires its number of dark channels in successive samples, and a single
        sample of a line in between resets the count;
  \item the steering turns towards the line, saturates at its limit, and its derivative
        term opposes an error that is closing;
  \item the station file holds three bays and three drop zones, the table matches colours
        to zones, and an unknown colour is rejected.
\end{itemize}
All 13 tests pass on the code as measured in \cref{sec:results}.

\subsection{Scoring Against Ground Truth}
A mission can report success and still have failed, for example when a coin slips out of
the jaws on the way and the robot releases nothing. The verdict on a run is therefore not
taken from the mission but from the referee of \cref{sec:referee}, which reads the true
positions of the coins. The two are separated in the software as well: the ground truth is
bridged from the simulator on a topic of its own, and no node of the robot subscribes to it.

\subsection{Recording}
The measurements of \cref{sec:results} were taken by a recording node that only listens. It
logs, at 10\unit{Hz} of simulated time, the true pose of the robot, the pose estimated by
the localisation and the pose of the odometry, together with every change of mission state,
every command to the jaws and every line of the referee. The scripts that run the missions
and evaluate the recordings are kept with the report, so that every number in it can be
reproduced.
